\section{Lasagna modules and the comparison framework}

We recall only the features needed for the comparison.  For a smooth
compact oriented four-manifold $W$ and a framed link $L\subset\partial W$,
the skein lasagna module $\SN(W;L)$ is generated by decorated surfaces in
$W$, modulo local relations coming from $\KhR_N$ cobordism maps
\cite{MorrisonWalkerWedrich2019,ManolescuWalkerWedrich2023}.

For a handle decomposition
\[
W_1\subset W_2\subset W_3\subset W_4,
\]
where $W_1$ is a one-handlebody, $W_2$ adds two-handles along a framed link,
$W_3$ adds three-handles along embedded spheres, and $W_4$ adds four-handles,
the published results have the following roles.
\begin{itemize}
\item MWW Theorem 3.2 identifies the two-handle stage with a cabled skein
lasagna quotient, whose relations include beta and psi operations.
\item MWW Theorems 3.7 and 3.10 describe each three-handle as a coequalizer
of the two hemisphere maps, and give the complete handle formula.
\item MWW Proposition 3.4 says that a four-handle attachment induces an
isomorphism for the empty boundary link.
\item MWW Corollary 3.5 gives
$\SN(S^4)\cong\Z$, concentrated in bidegree $(0,0)$.
\end{itemize}

At $N=2$, a nonzero class in bidegree $(0,494)$ would therefore obstruct a
diffeomorphism to $S^4$.  The intended proof constructs a functional at the
one-handle stage, shows that it annihilates the two-handle and three-handle
relations, and transports the resulting nonzero class through the
four-handle isomorphism.

The formal algebra correctly captures the last sentence.  The next sections
formulate the geometric identifications required for P0, C and S and record
which parts are established and which remain \Open.
